% Copyright 2026 The terCAD team. All rights reserved.
% Use of this content is governed by a CC BY-NC-ND 4.0 license that can be found in the LICENSE file.

\subsection{Emerging Design}

An entry point differs per each system:

\begin{lstlisting}
// Android, composeApp/src/androidMain/kotlin:
class MainActivity : ComponentActivity() {
  override fun onCreate(savedInstanceState: Bundle?) {
    super.onCreate(savedInstanceState)
    setContent { App() }
  }
}
// iOS, composeApp/src/iosMain/kotlin:
fun MainViewController() = ComposeUIViewController { App() }
// Desktop, composeApp/src/desktopMain/kotlin:
fun main() = application {
  Window(onCloseRequest = ::exitApplication) { App() }
}
// Web, composeApp/src/wasmJsMain/kotlin/main.kt:
fun main() {
  CanvasBasedWindow(canvasElementId = "ComposeTarget") { App() }
}
\end{lstlisting}

\noindent A core entry point is \q{shared/src/App.kt}-file for our application:

\begin{lstlisting}
@Composable
@Preview
@DevelopmentEntryPoint
fun App() {
  // Set the look of the application
    MaterialTheme {
        // To manage state (mutable -- change can be observed)
        var showContent by remember { mutableStateOf(false) }
        // Control the layout
        Column(
            modifier = Modifier
                .background(MaterialTheme.colorScheme.primaryContainer)
                .safeContentPadding()
                .fillMaxSize(),
            horizontalAlignment = Alignment.CenterHorizontally,
        ) {
            Button(onClick = { showContent = !showContent }) {
                Text("Click me!")
            }
            // Show or hide nested objects with an animation
            AnimatedVisibility(showContent) {
                val greeting = remember { Greeting().greet() }
                Column(
                    modifier = Modifier.fillMaxWidth(),
                    horizontalAlignment = Alignment.CenterHorizontally,
                ) {
                    Image(painterResource(Res.drawable.compose_multiplatform), null)
                    Text("Compose: $greeting")
                }
            }
        }
    }
}
\end{lstlisting}

It might be used \q{voyager}-router by adding its libraries into the project by updating 
\q{gradle/libs.versions.toml}-file in accordance with instructions from 
\href{https://voyager.adriel.cafe/setup}{https://voyager.adriel.cafe/setup}-page.
A Kotlin Multiplatform project can depend on multiplatform libraries that work for all target platforms:

\begin{lstlisting}
// ./libs.versions.toml
[versions]
voyager = "1.1.0-beta02"

[libraries]
voyager-navigator = { module = "cafe.adriel.voyager:voyager-navigator", version.ref = "voyager" }
voyager-tab-navigator = { module = "cafe.adriel.voyager:voyager-tab-navigator", version.ref = "voyager" }
voyager-transitions = { module = "cafe.adriel.voyager:voyager-transitions", version.ref = "voyager" }
\end{lstlisting}
